
\section{Fourier/Feynman representation of the Yukawa triple overlap}

We start from the screened Yukawa profile
\begin{equation}
G_m(r) \equiv \frac{e^{-mr}}{r},
\qquad
G_m(r)=\int \frac{d^3k}{(2\pi)^3}\,\frac{4\pi}{k^2+m^2}\,e^{i k\cdot r}.
\end{equation}

For three centers at positions $x_1,x_2,x_3$, define the exact overlap kernel
\begin{equation}
I_Y(x_1,x_2,x_3;m)
=
\int_{\mathbb R^3} d^3x\;
G_m(|x-x_1|)\,G_m(|x-x_2|)\,G_m(|x-x_3|).
\end{equation}

\subsection{Momentum-space representation}

Insert the Fourier representation for each Yukawa factor:
\begin{align}
I_Y
&=
\int d^3x
\prod_{a=1}^3
\left[
\int \frac{d^3k_a}{(2\pi)^3}
\frac{4\pi}{k_a^2+m^2}
e^{i k_a\cdot (x-x_a)}
\right].
\end{align}

The $x$-integration gives a momentum-conserving delta:
\begin{equation}
\int d^3x\,e^{i(k_1+k_2+k_3)\cdot x}
=
(2\pi)^3 \delta^{(3)}(k_1+k_2+k_3).
\end{equation}

Eliminate one momentum, e.g. $k_3=-(k+q)$, with $k\equiv k_1$, $q\equiv k_2$. Then
\begin{equation}
I_Y
=
\frac{(4\pi)^3}{(2\pi)^6}
\int d^3k\,d^3q\;
\frac{
e^{-i k\cdot r_{12}}\,e^{-i q\cdot r_{13}}
}{
(k^2+m^2)(q^2+m^2)\big((k+q)^2+m^2\big)
},
\end{equation}
where
\begin{equation}
r_{12}=x_1-x_2,\qquad
r_{13}=x_1-x_3.
\end{equation}

Thus the Yukawa overlap is exactly a three-propagator ``sunset''-type integral in momentum space.

\subsection{Feynman parameters}

Introduce
\begin{equation}
A=k^2+m^2,\qquad B=q^2+m^2,\qquad C=(k+q)^2+m^2.
\end{equation}
Use
\begin{equation}
\frac{1}{ABC}
=
2\int_0^1 d\alpha_1 d\alpha_2 d\alpha_3\,
\delta(1-\alpha_1-\alpha_2-\alpha_3)\,
\frac{1}{(\alpha_1 A+\alpha_2 B+\alpha_3 C)^3}.
\end{equation}

Then
\begin{equation}
\alpha_1 A+\alpha_2 B+\alpha_3 C
=
(\alpha_1+\alpha_3)k^2+(\alpha_2+\alpha_3)q^2+2\alpha_3 k\cdot q + m^2.
\end{equation}

Introduce a Schwinger parameter $s$:
\begin{equation}
\frac{1}{D^3}
=
\frac{1}{2}\int_0^\infty ds\, s^2 e^{-sD}.
\end{equation}
The factor $2$ from the Feynman parametrization cancels the factor $1/2$ above.

Thus
\begin{equation}
I_Y
=
\frac{(4\pi)^3}{(2\pi)^6}
\int_{\Delta_2} d\alpha\,
\int_0^\infty ds\, s^2 e^{-s m^2}
\int d^3k\,d^3q\,
\exp\left[
-s\Big((\alpha_1+\alpha_3)k^2+(\alpha_2+\alpha_3)q^2+2\alpha_3 k\cdot q\Big)
-i k\cdot r_{12}
-i q\cdot r_{13}
\right],
\end{equation}
where $\Delta_2$ is the simplex
\begin{equation}
\alpha_i\ge 0,\qquad \alpha_1+\alpha_2+\alpha_3=1.
\end{equation}

\subsection{Gaussian integration over momenta}

The quadratic form is controlled by the $2\times 2$ matrix
\begin{equation}
M=
\begin{pmatrix}
\alpha_1+\alpha_3 & \alpha_3\\
\alpha_3 & \alpha_2+\alpha_3
\end{pmatrix},
\qquad
\Delta \equiv \det M
=
\alpha_1\alpha_2+\alpha_1\alpha_3+\alpha_2\alpha_3.
\end{equation}

Performing the Gaussian integral over the six momentum components gives
\begin{equation}
\int d^3k\,d^3q\; e^{-s\,Q(k,q)-i k\cdot r_{12}-i q\cdot r_{13}}
=
\frac{\pi^3}{s^3 \Delta^{3/2}}
\exp\left[
-\frac{
(\alpha_2+\alpha_3) r_{12}^2 + (\alpha_1+\alpha_3) r_{13}^2 - 2\alpha_3\, r_{12}\cdot r_{13}
}{
4 s \Delta
}
\right].
\end{equation}

Now use
\begin{equation}
r_{12}\cdot r_{13}
=
\frac{d_{12}^2+d_{13}^2-d_{23}^2}{2},
\end{equation}
so that the exponent simplifies to
\begin{equation}
\frac{
\alpha_2 d_{12}^2+\alpha_1 d_{13}^2+\alpha_3 d_{23}^2
}{
4 s \Delta
}.
\end{equation}

All prefactors collapse:
\begin{equation}
\frac{(4\pi)^3}{(2\pi)^6}\,\pi^3 = 1.
\end{equation}

Therefore
\begin{equation}
I_Y
=
\int_{\Delta_2} d\alpha\;
\Delta^{-3/2}
\int_0^\infty \frac{ds}{s}\,
\exp\left[
-s m^2
-
\frac{
\alpha_2 d_{12}^2+\alpha_1 d_{13}^2+\alpha_3 d_{23}^2
}{
4 s \Delta
}
\right].
\end{equation}

\subsection{Bessel reduction}

Use the standard identity
\begin{equation}
\int_0^\infty \frac{ds}{s}\,e^{-\beta s-\gamma/s}
=
2 K_0\!\left(2\sqrt{\beta\gamma}\right),
\qquad
\beta>0,\ \gamma>0.
\end{equation}

With
\begin{equation}
\beta = m^2,
\qquad
\gamma=
\frac{
\alpha_2 d_{12}^2+\alpha_1 d_{13}^2+\alpha_3 d_{23}^2
}{
4\Delta
},
\end{equation}
we obtain the exact Fourier/Feynman form
\begin{equation}
\boxed{
I_Y(d_{12},d_{13},d_{23};m)
=
2\int_{\Delta_2} d\alpha\;
\Delta^{-3/2}\,
K_0\!\left(
m\sqrt{
\frac{
\alpha_2 d_{12}^2+\alpha_1 d_{13}^2+\alpha_3 d_{23}^2
}{
\Delta
}
}
\right)
}
\end{equation}
with
\begin{equation}
\Delta=
\alpha_1\alpha_2+\alpha_1\alpha_3+\alpha_2\alpha_3,
\qquad
\alpha_3=1-\alpha_1-\alpha_2.
\end{equation}

\subsection{Interpretation}

This shows:

\begin{enumerate}
\item The Yukawa triple overlap is not an arbitrary 3D spatial integral; it is a well-structured three-propagator integral.
\item The natural variables are
\begin{equation}
t = m\,s_{\max},
\qquad
q_1=\frac{s_{\min}}{s_{\max}},
\qquad
q_2=\frac{s_{\mathrm{mid}}}{s_{\max}},
\end{equation}
so that
\begin{equation}
I_Y = F(t;q_1,q_2).
\end{equation}
\item Fourier space does \emph{not} immediately produce an elementary algebraic closed form, but it does produce a clean propagator representation and a sharper starting point for further reduction.
\end{enumerate}

\subsection{Next possible directions}

The Fourier/Feynman form suggests three natural next steps:

\begin{enumerate}
\item search for a better shape-space parametrization (Jacobi / hyperradius / hyperangles),
\item derive asymptotic expansions directly from the Bessel kernel,
\item derive a differential equation for $I_Y$ as a function of triangle geometry.
\end{enumerate}
